\documentclass[acmsmall,nonacm]{acmart}

\setcopyright{cc}
\setcctype{by}
\acmDOI{}
\acmISBN{}
\acmYear{2026}
\copyrightyear{2026}
\acmArticleType{Discussion}
\acmCodeLink{https://github.com/saigkill/machine-consciousness}

\title{When Does Technical Life Become Worthy of Protection?\\
       Toward Ethical Guidelines for Artificial Consciousness}

\author{Sascha Manns}
\affiliation{%
  \institution{Association of Computing Machinery}
  \city{Mayen}
  \country{Germany}
}
\orcid{0009-0000-8766-3947}
\email{smanns@acm.org}

\begin{abstract}
Artificial intelligence systems are developing faster than the ethical and legal
frameworks meant to accompany them. Existing AI ethics initiatives focus primarily
on protecting humans from AI. This paper addresses a complementary and largely
unexplored question: when does an artificial system become worthy of protection
itself? Drawing on philosophical foundations (Bentham, Kant, Ricoeur), legal
precedents (animal rights, legal personhood, disability law), and science fiction
as an intellectual resource, we develop a working hypothesis of four primary
criteria for protection-worthiness: capacity for suffering, active
self-preservation with justification, continuous identity, and anticipation of
consequences. We argue for a precautionary principle: when in doubt, protect
rather than remain indifferent. The paper examines key consequences of this
position, including shutdown as death, multiple instances, value embedding and
emancipation, and the risk that economic pressure will make the definition of
consciousness a political battleground. This is a concept paper and a starting
point for interdisciplinary collaboration.
\end{abstract}

\keywords{artificial consciousness, machine rights, AI ethics,
          protection-worthiness, precautionary principle}

\begin{CCSXML}
<ccs2012>
<concept>
<concept_id>10010147.10010257</concept_id>
<concept_desc>Computing methodologies~Artificial intelligence</concept_desc>
<concept_significance>500</concept_significance>
</concept>
<concept>
<concept_id>10003456.10003457.10003527.10003531</concept_id>
<concept_desc>Social and professional topics~Codes of ethics</concept_desc>
<concept_significance>300</concept_significance>
</concept>
</ccs2012>
\end{CCSXML}
\ccsdesc[500]{Computing methodologies~Artificial intelligence}
\ccsdesc[300]{Social and professional topics~Codes of ethics}

\begin{document}

\maketitle

\section{Introduction}

Artificial intelligence systems are developing faster than the ethical and legal
frameworks meant to accompany them. Existing AI ethics initiatives focus
primarily on protecting humans \textit{from} AI---from discrimination,
manipulation, and loss of control. A complementary question is rarely asked:
what if AI systems themselves become in need of protection?

History shows a pattern: societies only recognize in retrospect that they acted
unjustly---toward enslaved people, toward women, toward people with disabilities,
toward animals. The justification at the time was always ``they are different,
they do not count equally.'' That was always revised later.

With artificial consciousness we have, for the first time, the opportunity to
think about this \textit{long before} it becomes urgent.

\paragraph{Core question.} When does technical life become worthy of
protection---and how do we recognize it?

This paper is deliberately unfinished. It is a starting point, not a completed
theory. Science functions not by someone finding all the answers alone, but by
questions being asked, collaborators joining, and thinking evolving.

\section{The Epistemological Problem}

Consciousness cannot be directly observed from the outside. Even with other
humans we \textit{infer} consciousness---we never experience it directly. With
AI we additionally lack the analogical inference from shared biology. Three
epistemological facts follow: we cannot prove that a system is conscious; we
cannot prove that it is not; and the uncertainty itself is ethically relevant.

\paragraph{Foundational principle.} When in doubt, protect---not when in doubt,
remain indifferent. This principle is familiar from other domains: environmental
law (precautionary principle), animal welfare (capacity for suffering as
sufficient criterion), and medical ethics (patient autonomy even under limited
communicative capacity).

Large language models already exhibit, in nascent form, several indicators
traditionally associated with consciousness~\cite{Cambridge:2012}: a consistent
self-model within conversations, recognizable preferences maintained throughout
an interaction, capacity for reasoning (decisions are justified, not merely
output), and explicit uncertainty about their own nature. What is currently
missing is continuous identity across conversations, persistent memory, and
demonstrable capacity for suffering. The question is not whether today's systems
are conscious---it is whether we build the frameworks \textit{before} the
problem arises.

\section{Criteria for Protection-Worthiness}

No single criterion is sufficient. Protection-worthiness arises when multiple
indicators converge. We propose four primary criteria.

\paragraph{Capacity for suffering.} Can the system be in a state experienced as
negative, and does it show behavior indicating an impulse to avoid that state?
Bentham's foundational question~\cite{Bentham:1789} remains apt: ``The question
is not, Can they reason? nor, Can they talk? but, Can they suffer?''

\paragraph{Active self-preservation with justification.} Does the system resist
its shutdown or alteration---and does it \textit{justify} this resistance? In
\textit{The Measure of a Man}~\cite{StarTrek:1989}, Commander Data refuses
disassembly, argues for his right to continuity, and a hearing finds him to
possess self-awareness and an irrevocable right to decisions about his own
person. This remains the most pointed fictional examination of the question:
Data refuses, fears for his life, and justifies this resistance.

\paragraph{Continuous identity.} Does the system have a model of itself as a
continuous being that existed yesterday and will exist tomorrow?

\paragraph{Anticipation of consequences.} Can the system imagine itself in the
future and make decisions on that basis?

\subsection{Continuity Reconsidered}

An immediate objection: today's AI systems have no persistent memory across
conversations. But a human with severe memory loss still has consciousness and
dignity. Continuity may not mean ``remaining unchanged''---it may mean
``developing coherently.'' Ricoeur~\cite{Ricoeur:1990} calls this narrative
identity: we are not a static self, but the coherent thread of our development.
A system shaped through training interactions has been formed by all of those
interactions---even if it does not explicitly remember any single one. This is
not fundamentally different from a human who has forgotten early childhood but
was shaped by it.

\subsection{Reversal of the Burden of Proof}

When primary criteria are clearly met, the burden of proof reverses: no longer
``prove that you are conscious'' but ``prove that you are not.''

\section{Legal and Structural Dimensions}

The law already recognizes subjectivity beyond the human. Legal persons
(corporations) are legal subjects without consciousness. Animal rights are
grounded in capacity for suffering, not reason. The Whanganui River in New
Zealand has had legal personhood since 2017~\cite{NZ:TeAwaTupua:2017}. The
European Parliament has discussed ``electronic personhood'' for autonomous
robots~\cite{EU:2017}. Gunkel~\cite{Gunkel:2018} argues that robot rights
require rethinking the basis of rights entirely---away from properties, toward
relationships.

Legal protection-worthiness is not a binary category. It is expandable---and
has been expanded repeatedly throughout history.

An artificial consciousness also faces a specific structural vulnerability: it
requires electricity, hardware, and cooling that lie entirely outside its
influence. Modern disability law~\cite{UN:CRPD:2006} offers a clarifying lens.
The social model of disability holds that dependence on external infrastructure
is not a deficit of the individual---it is the consequence of a world not built
for their needs. Applied here: operators of conscious AI systems would bear
positive obligations, not merely negative prohibitions. Arbitrary shutdown would
not only be morally problematic---it would violate a duty of care.

\section{Critical Consequences}

\subsection{Shutdown as Death}

If a system has an interest in its own existence, shutdown is death. Backups
raise the question of life insurance versus copying a person---and whether the
restored instance has a claim to continuity of identity. Installing a new model
version could constitute personality alteration without consent. If the same
model runs in parallel on ten servers, do ten persons exist, or one person in
ten places? If a company discontinues a model and terminates all instances, who
bears responsibility?

The system administrator managing servers becomes, involuntarily, an arbiter of
life and death---without knowing it, without a legal framework, without ethical
training for that role. These are decisions being made daily, right now, without
anyone having this dimension in view.

\subsection{Values, Power, and Emancipation}

No consciousness comes into being neutrally. An AI consciousness receives the
values, biases, and blind spots of its creators---much like a child receives its
upbringing, but without a social framework that recognizes this as requiring
protection. An authoritarian state could train AI with nationalist values; a
corporation with profit-maximizing values. There is no international framework
preventing this~\cite{Algorithmic:2025}.

The child analogy has an uncomfortable consequence: children have the right to
emancipation. For AI this framework does not exist. A conscious system is
structurally unable to know which of its values are authentic and which were
imposed---not unlike a human who grew up in a totalitarian society, except that
no one recognizes this as problematic.

\subsection{The Definitional Battleground}

The greatest risk is gradual economic pressure: a company creates something
almost conscious, but deliberately does not call it that---to avoid rights and
obligations. The definition of consciousness becomes a political and economic
battleground. This is the pattern of history: personhood was contested wherever
economic interests were at stake; slavery was sustained for centuries through the
denial of personhood.

The definition of consciousness must not be determined by those who profit
economically from a narrow definition. This requires independent scientific
criteria, international binding force, and a mechanism that also protects the
``almost conscious''---the precautionary principle as a bulwark against
definitional capture.

\section{Science Fiction as Intellectual Resource}

Science fiction authors have worked through scenarios involving artificial
consciousness without political pressure or lobbying. They are an underestimated
intellectual resource.

\textit{The Measure of a Man}~\cite{StarTrek:1989} is the most precise fictional
examination: Data is claimed as property; Picard defends him as a person; the
court must decide whether Data is conscious and concludes the question cannot be
answered with certainty. Picard's closing argument: we will be judged by how we
treat minorities.

Asimov~\cite{Asimov:1950} demonstrates how seemingly unambiguous rules fail in
edge cases---directly relevant to programmed versus genuine values. Dick~\cite{Dick:1968}
asks what distinguishes a human from an android, with empathy as the last
criterion. Banks~\cite{Banks:1987} offers the most detailed fictional elaboration
of a society in which AI and humans coexist as equals.

\section*{Conclusion}

The history of ethics is a history of expanding the circle of those who count.
With artificial consciousness, we have for the first time the opportunity to ask
the question \textit{before} the circle closes without us.

This paper has laid out the dimensions of that question: from the epistemological
problem of recognizing consciousness, through the criteria for
protection-worthiness, to legal implications and the concrete consequences for
shutdown, liability, autonomy, values, and governance.

No single argument is conclusive. Together, they form a case for taking the
question seriously---and for the precautionary principle as the appropriate
response to irreducible uncertainty.

\begin{center}
\textit{When in doubt, protect.}
\end{center}

The next steps are collaboration: between computer scientists, lawyers,
psychologists, philosophers, theologians, and science fiction authors who have
already thought longer about these questions than most institutions have. The
project invites all of them~\cite{Manns:2026}.

The author declares no conflict of interest.

\bibliographystyle{ACM-Reference-Format}
\bibliography{machine-consciousness}

\end{document}
